%%%%%%%%%%%%%%%%%%%%%%%%%%%%%%%%%%%%%%%%%%%%%%%%%%%%%%%%%%%%%%%
%
% Welcome to Overleaf --- just edit your LaTeX on the left,
% and we'll compile it for you on the right. If you open the
% 'Share' menu, you can invite other users to edit at the same
% time. See www.overleaf.com/learn for more info. Enjoy!
%
%%%%%%%%%%%%%%%%%%%%%%%%%%%%%%%%%%%%%%%%%%%%%%%%%%%%%%%%%%%%%%%
\documentclass[10pt]{article}
\usepackage[a4paper, total={8in, 11in} ]{geometry}
\usepackage{graphicx}
\usepackage{amssymb}

\renewcommand{\baselinestretch}{1}

\title{}

\begin{document}
First document. This is a simple example, with no 
extra parameters or packages included. 
\par

\fontsize{12}{10}
\textbf{\underline{Definition 1.1:}}
Let A and B be two sets. We say that A is a subset of B, or that A is contained
in B if and only if each x $ \in $A is also an element of B. We write  A $ \subseteq $ B to say that “A is a
subset of B.” If there is an element of A not in B, we write A $ \nsubseteq $ B.
\par

\fontsize{12}{10}
\textbf{\underline{Definition 1.2:}}
The empty set, written $ \emptyset $, is the set that has no elements.
\par

\fontsize{12}{10}
\textbf{\underline{Definition 1.3:}}
If A and B are statements, their conjunction is the statement “A and B”.
A $ \wedge $ B is true exactly when both A is true and B is true
\par
hello this shit ass
\end{document}